\section{The rules dialect in full}
\label{app:dialect}

The dialect is unchanged from the v0.4 and v0.5 studies, so this appendix is a
summary and a delta rather than a re-derivation. Table~\ref{tab:dialect} is the
index: it lists every consequential choice against the engine switch that
changes it. The v0.4 report's Appendix A states each choice at the length
required to re-implement it, including the ask-legality predicate, the deal and
shuffle, declaration resolution, the arbitration orders and the scoring rule
\cite{fishbot04}; that treatment is accurate except on the two clauses of
\S\ref{app:dialect-corrections}. The two files that carry the definitions are
\filepath{engine/src/fish.hpp} (deck, indexing, the \texttt{Rules} record, the
legality predicate) and \filepath{engine/src/game.hpp} (the driver: declaration
rounds, turn handling, the forced endgame, the action cap).

\begin{table}[ht]
\centering
\caption{The rules dialect used for every experiment in this report unless
stated otherwise, and the engine switch that changes each choice. The switches
are command-line flags of the \texttt{fish} driver; the underlying fields are
members of \texttt{Rules} in \texttt{engine/src/fish.hpp}. Rows marked
``modelling choice'' are not fixed by any published statement of the rules.
Unchanged from the v0.4 and v0.5 studies.}
\label{tab:dialect}
\small
\begin{tabular}{@{}p{0.29\linewidth}p{0.46\linewidth}p{0.19\linewidth}@{}}
\toprule
Choice & Value used & Switch \\
\midrule
Deck composition & 54 cards in nine half-suits (low and high band of each suit,
plus the four eights and two jokers); nine cards per player &
\texttt{-{}-sets=8} gives the 48-card, eight-half-suit form \\
Ask legality & Live half-suit; asker holds another card of it and not the named
card; target is an opponent and still holds cards & --- \\
Turn transfer & A successful ask retains the turn; a miss passes it to the
target & --- \\
Out-of-turn declarations & Permitted at any moment, including during an
opponent's turn & \texttt{outOfTurnDeclare} (\texttt{-{}-no-out-of-turn}) \\
Cardless players may declare & Permitted &
\texttt{cardlessMayDeclare} (\texttt{-{}-no-cardless-declare}) \\
Incorrect declaration & Awards the half-suit to the opposing team & --- \\
Declaration aftermath & The six cards leave play and the turn returns to the
player who held it & --- \\
Cardless turn-holder & Chooses which live teammate receives the turn, from its
own belief (a modelling choice, and more restrictive than the rules of record;
\S\ref{sec:rules-dialect}) & --- \\
Forced endgame (whole team cardless) & The live team declares every remaining
half-suit, sharing only willingness, over an eight-level ladder
$\{0.995, 0.98, 0.95, 0.90, 0.80, 0.65, 0.50, \text{best guess}\}$ &
\texttt{-{}-legacy} gives the two-level ladder $\{0.38, \text{best guess}\}$ \\
Simultaneous declarations & Lowest seat index declares first (a modelling
choice; cost measured in \S\ref{sec:rules-dialect}) &
\texttt{declArbitration} (\texttt{-{}-arb=low\textbar high\textbar turn}) \\
Action cap & 400 asks (360 under the v0.3 dialect), any residue adjudicated
neutrally by majority physical holding, ties to the holder of the lowest card &
\texttt{-{}-maxasks} \\
Conventions & No prearranged signalling conventions; team communication is
limited to the willingness bit described above & --- \\
\bottomrule
\end{tabular}
\end{table}

\subsection{Clause table and what changed}
\label{app:dialect-table}

Nothing in Table~\ref{tab:dialect} changed between \vfour{}, \vfive{} and
\vsix{}. The engine repairs of \S\ref{sec:engine-repairs} are inference and
bookkeeping defects, not rule changes, and none of them alters a game's outcome
given the same actions. Two entries in the table are worth reading against the
v0.4 appendix rather than in place of it. The cardless-turn-holder row is
narrower than the rules of record permit, and the row for simultaneous
declarations is a choice the rules do not make at all; both are quantified in
\S\ref{sec:rules-dialect}, and both measure at approximately zero, which is why
neither is a \vsix{} mechanism.

Dialect robustness is reported rather than assumed: \S\ref{sec:results-dialects}
replays the primary head-to-head result under \vsixDialectN{} perturbations of
this table over \vsixDialectDeals{} deals at seed \vsixDialectSeed{}, spanning
\vsixDialectSpan{}\%.

\subsection{Clauses the \vsix{} mechanisms depend on}
\label{app:dialect-depends}

Four clauses carry weight in this report, in the sense that a mechanism of
\S\ref{sec:mechanisms} or a measurement of \S\ref{sec:belief} would be different
under a different choice. They are stated here in full so that the dependence is
auditable.

\paragraph{Residue is adjudicated by physical majority.} The driver counts asks
and stops the game at \texttt{maxAsks}, 400 under the dialect studied. Awarding
the unresolved residue to either side would bias the outcome towards whichever
policy stalls more readily, so each remaining half-suit is instead awarded to
the team physically holding the majority of its six cards, ties going to the
team holding the lowest-indexed card of that half-suit. Two consequences matter.
An unclaimed half-suit is worth its majority holding rather than nothing, which
is the premise \S\ref{app:dialect-corrections} corrects; and the cap is a
backstop against unbounded play rather than part of the rules, so termination is
an empirical property of the fitted policies and is reported as one, at
\vsixLimitRate{}\% incidence.

\paragraph{The cardless-player rule.} A seat with no cards leaves the asking
cycle: it cannot ask, and it is an illegal target, because ask legality requires
the target to hold cards. It remains in the game in two ways. It can be named in
a teammate's declaration as the holder of a card --- impossible while it is
genuinely empty, so its emptiness constrains the allocations a teammate may
name, which is a hard constraint the exact count law of
\S\ref{sec:belief-count} consumes. And, with \texttt{cardlessMayDeclare}
enabled, it may itself declare. The turn can reach a cardless seat only
immediately after a declaration removed its last cards, and that seat then
chooses which live teammate receives it.

\paragraph{An opponent void is permanent.} Cards move in exactly one direction:
from a target to an asker, and only on a successful ask. Asking in half-suit $h$
requires the asker to hold another card of $h$. A seat holding no card of $h$
therefore can neither ask in $h$ nor receive a card of $h$, and no sequence of
legal actions restores it. Voids are monotone, they are public, and a void
created against an opponent removes that opponent's right to ask in the
half-suit for the rest of the game. This is a rules-level fact rather than a
tactical one, and it is what the offensive-void term of \S\ref{sec:mech-cheap}
prices.

\paragraph{The forced endgame.} When one team holds no cards at all, no legal
ask remains for either side and the live team must declare every remaining
half-suit. Teammates may negotiate openly but may exchange nothing except
willingness to declare, so the driver implements the selection as a threshold
sweep over the eight-level ladder of Table~\ref{tab:dialect} rather than as a
comparison of confidences: for each threshold in descending order the driver
scans the live team's seats over the remaining half-suits and asks each whether
it is willing to declare at confidence at least that threshold; the first seat
that says yes declares, the ladder restarts, and the sweep repeats until every
half-suit is resolved. Declarations here resolve exactly as elsewhere, including
the possibility of misdeclaring and handing a half-suit to the cardless team,
and are recorded separately in the event history. If every seat refuses at every
level, the remainder is awarded by physical majority as above.

\subsection{Two clauses corrected}
\label{app:dialect-corrections}

The v0.4 appendix is accurate on every clause but two. Both errors are stated in
that report and repeated in the source; neither was retracted by the v0.5 paper,
which measured them and did not print the result. They are restated correctly
here so that a reader who reaches this appendix first is not misled, and
\S\ref{sec:corr-dialect} gives them as corrections with the file and line of the
text being corrected.

\paragraph{An undeclared half-suit does not score nothing.} The v0.4 report's
declaration section justifies the unconditional second stage of its forcing rule
--- declare the best candidate whatever its estimated probability, past a second
horizon --- on the ground that an undeclared half-suit scores nothing. Inside
this harness it does not. The action cap calls the residue adjudicator described
above, which awards each unresolved half-suit by physical majority, and on a
wrong post-horizon declaration the declaring team holds a mean
\numResidueHeld{} of the six cards, which is the majority, which is what
adjudication would have given it. The stage does not convert nothing into a coin
flip; it converts a probable award into a certain concession. The premise is
true only of a dialect that awards nothing at the cap, and the v0.4 study
changed the harness away from that dialect after its parameter vector was
frozen, without updating the justification. The v0.4 rules section states the
dialect correctly, so that report contradicts itself within two sections.

\paragraph{Restarting the ladder sharpens nothing in observed play.} The v0.4
appendix states that a consequence of restarting the willingness ladder is that
early, safer declarations resolve cards and thereby sharpen the allocations
available for the later ones. The reasoning is correct and the situation does
not arise. Over every occasion on which a team went cardless with live
half-suits remaining --- \numLadderOccasions{} of them at seed
\numLadderSeed{} --- the number of live half-suits at that moment was
\numLadderLiveSets{} in every case, so there is nothing to order and every
ordering counterfactual is identical to the non-ordering one by construction.
The structure is not impossible: a cardless team requires only that the other
team hold a multiple of six cards. It does not occur because these policies cash
locked half-suits promptly. A policy made more patient about declaring would
begin to produce multi-half-suit forced endgames, at which point the restart
would start doing the work the v0.4 appendix credits it with, and the machinery
should be built alongside the patience change rather than after it.
